Short-range runoff prediction is important for hydrologic operations, yet the practical value of machine-learning models at one-day lead time depends strongly on benchmark design, information availability, and baseline strength. This study presents a chronology-respecting retrospective hindcast benchmark for next-day runoff prediction across a 50-basin CAMELS-US subset using open data only. Daily Daymet basin-mean meteorological forcing, USGS discharge, and static catchment attributes were assembled for 1990--2014, and next-day runoff was predicted from lagged runoff, rolling hydroclimatic summaries, seasonal terms, and basin descriptors. Four models were compared under a fixed temporal split: persistence, Ridge regression, Random Forest, and HistGradientBoosting. The benchmark is explicitly retrospective: it assumes end-of-day availability of observed same-day hydroclimatic information and therefore should not be interpreted as a real-time operational forecast system. Random Forest produced the highest pooled test skill, with NSE = 0.670, KGE = 0.756, RMSE = 2.177 mm d$^{-1}$, and MAE = 0.569 mm d$^{-1}$. Basin-wise summaries were consistent with the pooled ranking: Random Forest achieved the highest median basin-wise NSE (0.623) and was the best-performing model in 27 of 50 basins, while HistGradientBoosting ranked a close second. Persistence remained competitive, especially in KGE, indicating that short-memory hydrologic persistence still explains a large share of one-day-ahead predictability. The results therefore support a restrained conclusion: nonlinear tree ensembles provide credible improvement over strong simple baselines in this hindcast benchmark, but the gains are moderate rather than transformative. The manuscript emphasizes reproducibility, benchmark scope, and inference limits, and frames the contribution as benchmark design rather than algorithmic novelty.
